%============================================================================
% CAPÍTULO 5 — EXPERIMENTOS INICIAIS E RESULTADOS PRELIMINARES
% Estrutura evolutiva: preencher apenas o que tiver sido efetivamente executado.
% Não apresentar configurações planejadas como resultados obtidos.
%============================================================================
\chapter{Experimentos Iniciais e Resultados Preliminares}\label{cap:experimentos-iniciais}

Este capítulo apresenta o planejamento experimental e os resultados preliminares
obtidos com o método \textit{EFS4AC}, descrito no
Capítulo~\ref{cap:Metodologia}. O objetivo desta etapa da pesquisa é verificar
a viabilidade de sua integração com a classificação associativa e iniciar a
investigação dos efeitos das decisões metodológicas propostas sobre o desempenho
preditivo, a redução dimensional e a complexidade dos classificadores induzidos. Assim, o protocolo descrito
neste capítulo contempla as configurações efetivamente avaliadas nos experimentos iniciais.

%============================================================================
\section{Protocolo Experimental}\label{sec:exp-protocolo}
%============================================================================

Esta seção descreve os conjuntos de dados, os procedimentos de preparação dos dados,
o classificador associativo, as configurações do \textit{EFS4AC} e o procedimento
de avaliação.

%%%%%%%%%%%
\subsection{Conjuntos de Dados}\label{subsec:exp-bases}

A avaliação foi conduzida sobre conjuntos de dados de classificação com
atributo de classe nominal...COMPLETAR...

...VIDE TRABALHO MAICON...

%%%%%%%%%%%
\subsection{Pré-processamento}\label{subsec:exp-preprocessamento}

...VIDE TRABALHO MAICON...

%%%%%%%%%%%
\subsection{Classificador Associativo}
\label{subsec:exp-integracao}

...Identificar o classificador associativo e sua implementação; informar suporte, confiança, critérios de poda/ordenação e demais parâmetros...

%%%%%%%%%%%
\subsection{Configurações do EFS4AC e Métodos de Referência}\label{subsec:exp-configuracoes}

A arquitetura do \textit{EFS4AC} permite investigar quatro fatores principais:
a composição dos métodos-base, o controle da redundância entre \textit{rankings},
o operador de agregação e a regra de definição do subconjunto final. As
configurações previstas e avaliadas até este momento são sintetizadas na Tabela~\ref{tab:exp-configuracoes}.

%%%%%%%%%%%
\begin{table}[htbp]
    \centering
    \small
    \caption{Configurações metodológicas investigadas inicialmente.}
    \label{tab:exp-configuracoes}
    \renewcommand{\arraystretch}{1.15}
    \begin{tabular}{p{3.2cm}p{5.3cm}p{4.2cm}}
        \hline
        \textbf{Fator} & \textbf{Configurações avaliadas} &
        \textbf{Aspecto investigado} \\
        \hline
        Composição & $E_{U+M}$ & Efeito dos métodos-base \\
        Redundância & $D_0$, $D_L$ & Efeito do controle de similaridade \\
        Agregação & Borda Média, Mediana, Geométrica, L2; TOPSIS &
        Efeito do operador de consenso \\
        Corte & $K_{25}$, $K_{50}$, $K_{75}$, $K_{\Delta}$ &
        Efeito da quantidade selecionada \\
        Referências & Sem FS; métodos-base individuais &
        Comparação com alternativas sem consenso \\
        \hline
    \end{tabular}
\end{table}
%%%%%%%%%%%

As configurações de referência compreendem a indução do classificador com todos
os atributos (sem seleção) e a seleção realizada individualmente pelos métodos-
base. Essas comparações permitem investigar se os resultados do \textit{ensemble}
diferem daqueles obtidos por seus componentes isolados.

%%%%%%%%%%%
\subsection{Procedimento de Validação Cruzada e Prevenção de Vazamento de Informação}\label{subsec:exp-validacao}

Como os métodos-base utilizados pelo \textit{EFS4AC} empregam informações
da classe para estimar a relevância dos atributos, todo o processo de
seleção deve ser realizado exclusivamente com os dados de treinamento.
Isso inclui o pré-processamento, quando seus parâmetros
dependem dos dados, a construção dos \textit{rankings}, a análise da
diversidade, a agregação e a determinação do ponto de corte. Esse
procedimento busca evitar o vazamento de informações dos dados de teste
para as etapas de seleção de atributos e indução do classificador.

Em um procedimento de validação cruzada com $k$ partições
(\textit{folds}), seja $D_{\mathrm{tr}}^{(q)}$ o conjunto de treinamento
e $D_{\mathrm{te}}^{(q)}$ o conjunto de teste correspondente à
partição $q$, com $q=1,\ldots,k$. Em cada partição, o
\textit{EFS4AC} é executado novamente utilizando exclusivamente
os dados de treinamento, conforme a Equação~\ref{eq:exp-efs-fold}, em que $\operatorname{EFS4AC}(\cdot)$ representa a aplicação do método
proposto e $F_q^{*}$ corresponde ao subconjunto de atributos selecionado
a partir dos dados de treinamento da partição $q$.

%%%%%%%%%
\begin{equation}\label{eq:exp-efs-fold}
    F_q^{*} =
    \operatorname{EFS4AC}(D_{\mathrm{tr}}^{(q)})
\end{equation}
%%%%%%%%%

Após a seleção, os conjuntos de treinamento e teste são projetados
sobre os atributos pertencentes a $F_q^{*}$, conforme a
Equação~\ref{eq:exp-projecoes-fold}. Nessa equação, $\pi_{F_q^{*}}(\cdot)$ representa a operação de projeção
que preserva somente os atributos pertencentes ao subconjunto
$F_q^{*}$, além do atributo de classe.
Os conjuntos $D_{\mathrm{tr}}^{*(q)}$ e $D_{\mathrm{te}}^{*(q)}$
correspondem, respectivamente, aos dados de treinamento e teste
após a redução de atributos.

%%%%%%%%%
\begin{equation}\label{eq:exp-projecoes-fold}
    D_{\mathrm{tr}}^{*(q)}
    =
    \pi_{F_q^{*}}(D_{\mathrm{tr}}^{(q)}),
    \qquad
    D_{\mathrm{te}}^{*(q)}
    =
    \pi_{F_q^{*}}(D_{\mathrm{te}}^{(q)})
\end{equation}
%%%%%%%%%

O classificador associativo é então induzido utilizando
$D_{\mathrm{tr}}^{*(q)}$ e avaliado sobre
$D_{\mathrm{te}}^{*(q)}$. Esse procedimento é repetido para
todas as $k$ partições da validação cruzada. Neste trabalho, utilizou-se $k=10$.

Cabe destacar que diferentes partições podem produzir subconjuntos
$F_q^{*}$ distintos, uma vez que a seleção de atributos é realizada
independentemente em cada conjunto de treinamento. Dessa forma,
os dados de teste não participam da escolha dos atributos, permitindo
avaliar o desempenho do classificador associativo sem introduzir
vazamento de informações durante o processo de seleção.

%============================================================================
\section{Critérios de Avaliação}
\label{sec:exp-criterios}
%============================================================================

A avaliação do \textit{EFS4AC} considera não apenas o desempenho preditivo,
mas também o grau de redução dimensional, a estrutura do classificador
associativo, a estabilidade da seleção e o custo computacional. Essas dimensões
serão analisadas conjuntamente, sem pressupor que a melhora de uma medida
implique necessariamente melhora nas demais.

\textbf{Desempenho Preditivo.} A medida de desempenho será a F1-Macro, Equação~\ref{eq:exp-f1-macro}, em que $C$ representa o número de classes.
%%%%%%%%%%%
\begin{equation}\label{eq:exp-f1-macro}
    F1_{\mathrm{Macro}}=\frac{1}{C}\sum_{c=1}^{C}F1_c,
    \qquad
    F1_c=2*\frac{\operatorname{Precisao}_c*\,\operatorname{Revocacao}_c}
    {\operatorname{Precisao}_c+\operatorname{Revocacao}_c}
\end{equation}
%%%%%%%%%%%

\textbf{Redução Dimensional.} Para cada subconjunto $F^{*}$, serão registrados o número de atributos selecionados e a proporção de redução em relação ao conjunto original, conforme Equação~\ref{eq:exp-reducao}. Uma redução maior não será interpretada isoladamente como superior,
pois a exclusão de atributos informativos pode prejudicar o desempenho da classificação e a complexidade dp modelo.
%%%%%%%%%%%
\begin{equation}\label{eq:exp-reducao}
    \operatorname{Red}(F^{*})=
    \left(1-\frac{|F^{*}|}{|F|}\right)\times 100
\end{equation}
%%%%%%%%%%%

\textbf{Complexidade do classificador associativo.} Para analisar a estrutura do modelo, serão registrados inicialmente o número de regras finais e o tamanho médio dos antecedentes. Seja $\mathcal{Q}$ o conjunto de regras e $\operatorname{ant}(q)$ o antecedente da
regra $q$. Para $|\mathcal{Q}|>0$, define-se $\overline{L}_{\mathrm{ant}}$ de acordo com a Equação~\ref{eq:exp-tamanho-medio}.
%%%%%%%%%%%
\begin{equation}\label{eq:exp-tamanho-medio}
    \overline{L}_{\mathrm{ant}}=
    \frac{1}{|\mathcal{Q}|}\sum_{q\in\mathcal{Q}}|\operatorname{ant}(q)|
\end{equation}
%%%%%%%%%%%

\textbf{Estabilidade da Seleção.} A estabilidade será considerada como a consistência dos subconjuntos obtidos
em diferentes partições de treinamento. Para dois subconjuntos
$F_a^{*}$ e $F_b^{*}$, o índice de Jaccard é obtido via Equação~\ref{eq:exp-jaccard}. Para $k$ subconjuntos, a estabilidade média (Est) entre pares pode ser expressa via Equação~\ref{eq:exp-estabilidade}. Essa medida será interpretada em conjunto com o tamanho dos subconjuntos,
pois a sobreposição também é influenciada pela quantidade de atributos
selecionados.

%%%%%%%%%%%
\begin{equation}\label{eq:exp-jaccard}
    J(F_a^{*},F_b^{*})=
    \frac{|F_a^{*}\cap F_b^{*}|}{|F_a^{*}\cup F_b^{*}|}
\end{equation}
%%%%%%%%%%%

%%%%%%%%%%%
\begin{equation}\label{eq:exp-estabilidade}
    \operatorname{Est}=
    \frac{2}{k(k-1)}\sum_{a=1}^{k-1}\sum_{b=a+1}^{k}
    J(F_a^{*},F_b^{*})
\end{equation}
%%%%%%%%%%%

\textbf{Custo Computacional [OPCIONAL].} Serão registrados separadamente o tempo de seleção de atributos
$T_{\mathrm{EFS}}$ e o tempo de indução do classificador
$T_{\mathrm{CA}}$. Quando pertinente, será considerado o tempo total, i.e., $T_{\mathrm{total}}=T_{\mathrm{EFS}}+T_{\mathrm{CA}}$. Essa distinção permite verificar se o custo adicional da seleção é compensado
por eventual redução no custo de indução do classificador.

%===========================================================================
\section{Resultados e Discussão}
\label{sec:exp-resultados}
%============================================================================

%%%%%%%%%%%%%%%%%%%%%%
\subsection{Desempenho Preditivo e Redução Dimensional}\label{subsec:exp-resultados-f1-reducao}

Objetivo: verificar se a redução do conjunto de
atributos preserva, melhora ou prejudica o desempenho de classificação em
relação aos cenários de referência disponíveis.

%%%%%%%%%%%
\begin{table}[htbp]
    \centering
    \small
    \caption{Desempenho preditivo e redução dimensional: estrutura para os resultados preliminares.}
    \label{tab:exp-resultados-principais}
    \begin{tabular}{p{2.5cm}p{3.0cm}p{2.7cm}p{2.1cm}p{2.1cm}}
        \hline
        \textbf{Base} & \textbf{Configuração} &
        \textbf{F1-Macro} & \textbf{Atributos} & \textbf{Redução} \\
        \hline
        \multicolumn{5}{c}{\textit{Preencher com os resultados obtidos.}} \\
        \hline
    \end{tabular}
\end{table}
%%%%%%%%%%%

%%%%%%%%%%%%%%%%%%%%%%
\subsection{Efeitos sobre o Classificador Associativo}
\label{subsec:exp-resultados-regras}

Analisar: número de regras finais e o tamanho médio dos antecedentes. Reduções nessas medidas não serão interpretadas isoladamente como
melhoria, devendo ser consideradas em conjunto com a F1-Macro.

%%%%%%%%%%%%%%%%%%%%%%
\subsection{Estabilidade da Seleção}\label{subsec:exp-resultados-estabilidade}
...

%%%%%%%%%%%%%%%%%%%%%%
%Posteriormente - estudo de ablação
%\subsection{Componentes do EFS4AC}\label{subsec:exp-resultados-componentes}

%As saídas intermediárias do \textit{EFS4AC} também poderão ser utilizadas para compreender o comportamento das configurações avaliadas. A matriz de correlação de Spearman permitirá identificar concordâncias e redundâncias entre os métodos-base. A comparação dos \textit{rankings} consensuais e dos subconjuntos selecionados permitirá verificar como diferentes agregadores e regras de corte influenciam a seleção final.

%============================================================================%
\section{Considerações Finais}\label{sec:exp-consideracoes}

...FAZER...
